% ============================================================
% §2 · ¿Qué es el software y para qué sirve? — Manual de Usuario K-QGMIP
% ============================================================

\section{¿Qué es el software y para qué sirve?}

\textbf{K-QGMIP} es un \textit{framework} computacional en Python diseñado para calcular la
\textbf{k-Partición de Mínima Información} (\textit{k-MIP}) de sistemas causales discretos.
Dado un sistema de $n$ variables binarias y su Matriz de Probabilidad de Transición (TPM),
el \textit{framework} identifica la forma óptima de dividir el sistema en $k$ partes
independientes de modo que la pérdida de información integrada $\varphi$ sea mínima.
Este resultado revela la \textbf{estructura modular} del sistema: las divisiones que menos
perturban su dinámica causal son precisamente aquellas con menor información integrada.

\subsection*{¿Qué puede hacer el usuario con K-QGMIP?}

Con K-QGMIP el usuario puede:

\begin{itemize}
  \item \textbf{Seleccionar el número de particiones $k$}: analizar el sistema con
        $k \in \{2, 3, 4, 5\}$ partes, desde la bipartición clásica hasta divisiones
        en cinco subconjuntos independientes.

  \item \textbf{Elegir el método de cálculo}: el \textit{framework} ofrece dos estrategias
        complementarias que el usuario puede seleccionar según el tamaño del sistema y el
        tiempo de cómputo disponible:
        \begin{itemize}
          \item \textbf{KGeoMIP}: estrategia geométrica, adecuada para sistemas de tamaño
                moderado a grande. Reutiliza la tabla de costos de transiciones entre estados
                para evaluar candidatos de k-partición sin necesidad de búsqueda exhaustiva.
          \item \textbf{KQNodes}: estrategia submodular basada en el algoritmo de Queyranne,
                con menor complejidad asintótica para valores bajos de $k$.
                Prioritaria para sistemas con muchos nodos donde la eficiencia es crítica.
        \end{itemize}

  \item \textbf{Analizar sistemas de distintos tamaños}: desde redes pequeñas de validación
        (10--15 nodos) hasta sistemas de escala moderada con 20 o más nodos, cargando la TPM
        desde un archivo CSV estándar.

  \item \textbf{Obtener la partición óptima con su pérdida}: el resultado incluye la
        k-partición identificada, el valor de $\varphi$ asociado y metadatos del subsistema
        evaluado (estado inicial, máscaras binarias, tiempo de ejecución).

  \item \textbf{Ejecutar análisis en lote}: procesar múltiples subsistemas definidos en un
        archivo Excel, obteniendo resultados comparativos en formato CSV para análisis posterior.
\end{itemize}

\subsection*{Contexto de uso}

K-QGMIP está orientado a tres contextos principales:

\begin{enumerate}
  \item \textbf{Investigación académica en IIT 4.0}: el \textit{framework} permite estudiar
        cómo varía la estructura de mínima información de un sistema al aumentar el número de
        particiones permitidas $k$, aportando evidencia empírica sobre la organización modular
        de sistemas causales.

  \item \textbf{Experimentos exploratorios}: el usuario puede probar rápidamente distintas
        combinaciones de método, $k$ y subsistema para identificar configuraciones de interés
        antes de lanzar análisis sistemáticos de mayor costo computacional.

  \item \textbf{Análisis sistemáticos reproducibles}: la semilla aleatoria fija
        (\texttt{numpy seed = 73}) y la ejecución desde archivos de configuración garantizan
        que los resultados son deterministas y reproducibles entre sesiones y máquinas.
\end{enumerate}
